\section*{Chapitre \ref{ch:rc:stress-testing-and-scenarios} --- Tests de résistance et scénarios}

\begin{solution}{exo:rc:stress-testing-and-scenarios:1}
La \omterm{def:rc:market-risk-measures:var}{valeur en risque} (VaR) est un quantile d'une distribution estimée sur une fenêtre~; un scénario de stress est une seule trajectoire sévère
issue d'une autre période, au-delà du quantile et hors du régime de la fenêtre. Une VaR à 99~\% correcte
est dépassée, de quelque montant que ce soit, sur 1~\% des jours.
\end{solution}

\begin{solution}{exo:rc:stress-testing-and-scenarios:2}
$\E[x_2\mid x_1] = \rho\frac{\sigma_2}{\sigma_1}x_1 = 0.6\times\frac{5}{10}\times40 = 12$ points de base.
\end{solution}

\begin{solution}{exo:rc:stress-testing-and-scenarios:3}
$10/2.5 = 4$ écarts types.
\end{solution}

\begin{solution}{exo:rc:stress-testing-and-scenarios:4}
Le rendement à deux ans, corrélé à celui à dix ans, monte lui aussi (80 points de base), et le portefeuille est vendeur
d'obligations à deux ans~: cette jambe gagne et compense une partie de la perte sur le dix ans. Un choc sur un seul facteur
ignore les couvertures que procure le mouvement conjoint des autres \omterm{def:rc:market-risk-measures:factor}{facteurs de risque}.
\end{solution}

\begin{solution}{exo:rc:stress-testing-and-scenarios:5}
Le modèle linéaire ignore le gamma du straddle vendu~: son scénario le plus plausible répartit la perte entre les taux et les devises
à 3.93 écarts types et, en réévaluation complète, il perd 14.00~millions de dollars, soit plus que la cible. Le straddle perd
sur les fortes variations des devises dans un sens comme dans l'autre, si bien qu'un scénario plus proche,
à 2.47 écarts types, concentré sur l'EURUSD, atteint 10~millions de dollars.
\end{solution}

\begin{solution}{exo:rc:stress-testing-and-scenarios:6}
Que les mouvements seraient les mêmes sur les marchés actuels (niveaux, volatilités, corrélations, structure du marché),
que les positions actuelles soient conservées sur toute la période sans rééquilibrage,
et que le portefeuille pourrait être liquidé à ces prix.
\end{solution}

\begin{solution}{exo:rc:stress-testing-and-scenarios:7}
Sans le straddle, la distance linéaire est de 3.84 écarts types et celle en réévaluation complète de
3.95~: la convexité positive des obligations rend une perte légèrement plus difficile à atteindre que ne le dit le modèle linéaire,
l'inverse de l'effet du straddle.
\end{solution}

\begin{solution}{exo:rc:stress-testing-and-scenarios:8}
La probabilité repose sur une loi normale et sur la covariance récente~; les mouvements réels ont des queues épaisses et
les corrélations bondissent en période de crise, si bien que des mouvements de quatre sigmas sur dix jours surviennent bien plus souvent
(les épisodes historiques de ce chapitre coûtent deux à trois fois la VaR à dix jours). La distance classe
les scénarios~; ce n'est pas une probabilité.
\end{solution}

\begin{solution}{pb:rc:stress-testing-and-scenarios:1}
\textbf{1.} 1.5~milliard de dollars.
\textbf{2.} 905~millions de dollars.
\textbf{3.} Tant que le client paie, le courtier est couvert (il détient les actions en couverture des swaps)~; en cas de défaut,
les swaps disparaissent et le courtier reste long des actions, la marge étant son seul coussin.
\textbf{4.} Presque rien~: les swaps et les couvertures se compensent, et la marge couvre les mouvements quotidiens.
\textbf{5.} La position totale du client chez l'ensemble des courtiers~: les mêmes titres, en taille multipliée,
dont la vente forcée ferait baisser les cours au détriment de chaque courtier.
\textbf{6.} Baisses de 8.7, 7.9, 7.1, 6.3 et 5.4~\% pour les cinq actions, de la position la plus importante à la
plus petite.
\textbf{7.} 1.66 écart type~; une probabilité normale de 4.9~\% sur cinq jours.
\textbf{8.} Le scénario se charge sur la direction $\Sigma\delta$~: les actions à poids plus élevé contribuent davantage à
la perte par unité de mouvement, et le facteur commun les fait toutes bouger ensemble.
\textbf{9.} 7.5~\% sur chaque action.
\textbf{10.} La distance monte à 4.42 écarts types, soit une probabilité normale de $5.0\times10^{-6}$.
\textbf{11.} Les positions représentaient plusieurs jours de volume échangé~; le courtier est en concurrence avec d'autres courtiers
qui vendent les mêmes titres, et les cours sautent avant toute vente.
\textbf{12.} Ils ajoutent une baisse commune et corrélée de ces seules actions, que la covariance des jours normaux ne contient pas~:
le scénario est à la fois plus grand et plus probable.
\textbf{13.} Il avait été fixé bas et de manière statique, alors que la concentration et la volatilité des positions augmentaient~:
le coussin a diminué par rapport au risque.
\textbf{14.} La position totale du client (ou au moins sa concentration sur chaque titre chez ce courtier),
la liquidité exprimée en jours de volume, la volatilité, recalibrées chaque jour.
\textbf{15.} Une limite sur chaque titre exprimée en part de son volume quotidien et du capital du courtier,
la marge augmentant fortement avec la concentration.
\textbf{16.} La perte provient d'un événement conjoint (défaut et baisse concentrée) dans une position encombrée,
hors de toute fenêtre de jours normaux.
\textbf{17.} Comme déclencheur de marge, de limites et de sorties~: lorsque la distance passe sous un seuil, il faut agir.
\textbf{18.} Son exposition totale chez l'ensemble des courtiers, son effet de levier et sa liquidité.
\textbf{19.} Résultat nommé~: \emph{le family office vu du côté du courtier}~: des baisses de 8.7, 7.9, 7.1, 6.3 et
5.4~\% sur cinq jours épuisent la marge de 7.5~\%, à seulement 1.66 écart type.
\textbf{20.} Quels mondes plausibles feraient tomber la firme, et à quelle distance ils se trouvent.
\end{solution}

\begin{solution}{iq:rc:stress-testing-and-scenarios:1}
VaR~: une affirmation probabiliste sur les pertes ordinaires, estimée à partir de données récentes, utile pour les limites
au jour le jour. \omterm{def:rc:stress-testing-and-scenarios:stress}{Test de résistance}~: la perte dans un scénario sévère précisé, sans probabilité, utile pour les queues
de distribution, les concentrations et les changements de régime. Il faut utiliser les deux.
\iqlookfor{rôles complémentaires.}
\end{solution}

\begin{solution}{iq:rc:stress-testing-and-scenarios:2}
Préciser les chocs clés du scénario, compléter les autres facteurs par espérance conditionnelle sous une
covariance (éventuellement stressée) ou un \omterm{def:rc:structural-credit-models:structural}{modèle structurel}, vérifier les signes et les ordres de grandeur par rapport à l'historique,
puis faire relire le tout par les desks.
\iqlookfor{conditionnement et relecture experte.}
\end{solution}

\begin{solution}{iq:rc:stress-testing-and-scenarios:3}
Minimiser $\tfrac12x^\top\Sigma^{-1}x$ sous la contrainte $\delta^\top x = -L$~: $x^* = -L\Sigma\delta/(\delta^\top\Sigma\delta)$, distance $L/\sqrt{\delta^\top\Sigma\delta}$.
\iqlookfor{le lagrangien et le résultat.}
\end{solution}

\begin{solution}{iq:rc:stress-testing-and-scenarios:4}
Un choc instantané et sévère sur de nombreux facteurs de marché, appliqué aux positions de trading et à la juste valeur
des grandes banques, dont l'ampleur est fixée par le superviseur~; la VaR est une estimation statistique propre à la banque sur
un jour ou dix jours.
\iqlookfor{fixé par le superviseur, instantané, large.}
\end{solution}

\begin{solution}{iq:rc:stress-testing-and-scenarios:5}
Par client~: défaut plus une baisse concentrée de ses plus grandes positions sur un horizon de liquidation réaliste,
avec encombrement chez les différents courtiers~; \omterm{def:rc:stress-testing-and-scenarios:reverse}{test de résistance inversé} jusqu'à la marge~; limites sur la concentration et
sur le nombre de jours de volume~; marge dynamique.
\iqlookfor{conditionné au défaut, concentration et liquidité.}
\end{solution}

\begin{solution}{iq:rc:stress-testing-and-scenarios:6}
Épisodes historiques par classe d'actifs et par régime, \omterm{def:rc:stress-testing-and-scenarios:hypothetical}{scénarios hypothétiques} pour les risques actuels, résultats des \omterm{def:rc:stress-testing-and-scenarios:reverse}{tests de résistance inversés},
scénarios prudentiels~; définitions versionnées, responsables, cycles de revue, correspondances de facteurs pour les nouveaux produits,
et réconciliation des résultats avec la VaR et le P\&L.
\iqlookfor{contenu et gouvernance.}
\end{solution}
